% Optional source-only fragment for the round-four report.
\section{New helical and canonical-band clients}
\paragraph{H1: reuse integer-multiple periodicity.}
At the fixed ProveIt commit, \texttt{row\_zero\_of\_helix}
(\texttt{WangTiles942HelicalQuotient.lean}, lines 22--40) performs separate
positive- and negative-time induction.  Instead put
\(F(x,t)=T(x,t)\).  Its premise is exactly periodicity under \((w,-1)\).
The already available \texttt{Function.Periodic.zsmul} gives
\[
 F\bigl((x,t)+t\mathbin{\bullet}(w,-1)\bigr)=F(x,t),
 \qquad (x,t)+t\mathbin{\bullet}(w,-1)=(x+wt,0).
\]
Symmetry proves the original statement.  The pinned mathlib implementation
is \texttt{Mathlib/Algebra/Ring/Periodic.lean}, lines 155--158, at
\texttt{81a5d257c8e410db227a6665ed08f64fea08e997}.
No positivity condition is introduced: \(w=0\), negative widths and negative
times remain valid.  The candidate body shrinks from 16 to 7 physical lines.
This is a concrete reuse of an existing group-period theorem, not a claim
that a generic additive-action framework is novel.

\paragraph{H2: retain the concrete quotient before enriching its properties.}
The new \texttt{band\_quotient\_properties} in
\texttt{WangTiles942CycleBand.lean} proves, for the literal existing
\texttt{quotientWord (extend (band L H) ...)}, both cyclic matching and
\[
 \forall x\in[-L,R)\;\forall t\in\mathbb Z,\qquad
 \operatorname{tile}(a(\operatorname{address}(x,t)))
   =\operatorname{band}(L,H,x,t).
\]
The existing existential endpoint then forgets the witness expression;
the existing canonical endpoint combines the same certificate with its
separately established marker uniqueness.  The definition of
\texttt{canonicalWord} is unchanged, so its identity is preserved by
reduction rather than a new uniqueness premise.  The actual halt, all
independent margin/height lower bounds, signed time, the included left wall,
and the excluded right wall are retained.  No coprimality or square-band
restriction appears.

The two original endpoints are
\texttt{CycleBand.marked\_word\_of\_halt}, lines 136--165, and
\texttt{BandMarker.canonicalWord\_properties}, lines 79--106.
Their 40 proof-body lines become 21, including all 15 lines of the new
helper proof.  Charging its complete statement, documentation and spacing
leaves only a four-line whole-file reduction, not a nineteen-line net
source reduction.  Combined with H1, the three full modules change from
438 to 425 lines; selected proof bodies change from 56 to 28.  There is
one added theorem, no new definition, and no changed original signature.

\paragraph{Incremental novelty and validation.}
The final baseline audit found no corresponding concrete helical descent
or integer-period client among all 21 packet manuscripts and technical
notes.  Generic residue normalization, word wrappers, finite dispatch,
NorthView classification, translated intervals and prior group-action
suggestions remain assumed applied and are not counted here.  Ten
non-Lean tests pass, including source and mathlib blob identity, combined
patch replay, preserved signatures and word definition, signed transport,
and independently chosen dimension models.  Another researcher found no
semantic gap in source review.  All candidate Lean remains
\textbf{uncompiled}; there is no kernel, axiom-audit or performance claim.
